% Functional wiring from esp32.ino. It is not a manufactured carrier pinout.
\newcommand{\WearableWiringFigure}{%
\begin{figure}[H]\centering
\resizebox{\linewidth}{!}{%
\begin{tikzpicture}[font=\sffamily\small,
  sda/.style={draw=accent,line width=1pt},
  scl/.style={draw=blue!65!black,line width=1pt,dashed},
  supply/.style={draw=red!65!black,line width=1pt},
  ground/.style={draw=ink!80,line width=1pt},
  pinlabel/.style={font=\sffamily\scriptsize,anchor=west,fill=white,inner sep=1pt}]
\node[anchor=west,font=\sffamily\large\bfseries,text=ink] at (0,8.8) {Current wearable: shared bus, three sensors, one output};
\draw[fill=light,draw=ink!55,rounded corners=5pt] (0,-.2) rectangle (3.75,8.1);
\node[align=center,font=\sffamily\bfseries] at (1.875,7.55) {KEYES carrier\\+ classic ESP32};
\draw[draw=ink!50,fill=white,rounded corners=2pt] (.55,5.8) rectangle (3.15,6.8);
\node[font=\sffamily\scriptsize] at (1.85,6.3) {ESP32 acquisition / SPP};
\node[anchor=west,font=\sffamily\scriptsize,text width=2.9cm] at (.42,1.95) {USB supplies the development assembly.\\[4pt]Pin numbers name ESP32 GPIOs, not carrier header positions.};
\foreach \y/\label in {5.25/{GPIO21 / SDA},4.6/{GPIO22 / SCL},3.95/{module supply*},3.3/{GND},.8/{GPIO25 / output}} {
 \draw[fill=white,draw=ink!50] (3.64,\y-.08) rectangle (3.86,\y+.08);
 \node[anchor=east,font=\sffamily\scriptsize] at (3.5,\y) {\label};}
% Four named vertical nets; junction dots indicate genuine electrical joins.
\draw[sda] (3.86,5.25)--(4.7,5.25)--(4.7,7.7);\draw[sda] (4.7,5.25)--(4.7,-.6);
\draw[scl] (3.86,4.6)--(5.35,4.6)--(5.35,7.7);\draw[scl] (5.35,4.6)--(5.35,-.6);
\draw[supply] (3.86,3.95)--(6,3.95)--(6,7.7);\draw[supply] (6,3.95)--(6,-1.9);
\draw[ground] (3.86,3.3)--(6.65,3.3)--(6.65,7.7);\draw[ground] (6.65,3.3)--(6.65,-2.3);
\node[font=\sffamily\scriptsize,text=accent,rotate=90,anchor=west] at (4.58,7.65) {SDA};
\node[font=\sffamily\scriptsize,text=blue!65!black,rotate=90,anchor=west] at (5.23,7.65) {SCL};
\node[font=\sffamily\scriptsize,text=red!65!black,rotate=90,anchor=west] at (5.88,7.65) {VCC*};
\node[font=\sffamily\scriptsize,text=ink,rotate=90,anchor=west] at (6.53,7.65) {GND};
\foreach \base/\name/\addr/\detail in {
  6.6/{MAX30102}/{0x57}/{Red + infrared optical samples},
  3.75/{GY906 / MLX90614}/{0x5A}/{Object surface + sensor ambient},
  .9/{MPU6050}/{0x68}/{3-axis acceleration + 3-axis gyro}} {
 \draw[draw=ink!40,fill=white,rounded corners=3pt] (8.1,\base-1.02) rectangle (14.85,\base+1.16);
 \node[anchor=west,font=\sffamily\bfseries] at (9.25,\base+.72) {\name};
 \node[anchor=west,font=\sffamily\scriptsize,text=accent] at (9.25,\base+.25) {7-bit address \addr};
 \node[anchor=west,font=\sffamily\scriptsize,text width=5.1cm] at (9.25,\base-.35) {\detail};
 \draw[sda] (4.7,\base+.65)--(8.1,\base+.65);\fill[accent] (4.7,\base+.65) circle (.045);
 \draw[scl] (5.35,\base+.15)--(8.1,\base+.15);\fill[blue!65!black] (5.35,\base+.15) circle (.045);
 \draw[supply] (6,\base-.35)--(8.1,\base-.35);\fill[red!65!black] (6,\base-.35) circle (.045);
 \draw[ground] (6.65,\base-.85)--(8.1,\base-.85);\fill[ink] (6.65,\base-.85) circle (.045);
 \node[pinlabel] at (8.2,\base+.65) {SDA};\node[pinlabel] at (8.2,\base+.15) {SCL};
 \node[pinlabel] at (8.2,\base-.35) {VCC*};\node[pinlabel] at (8.2,\base-.85) {GND};
}
\node[archbox,text width=5.9cm,minimum height=1.5cm,anchor=west] (buzz) at (8.1,-1.95)
 {\textbf{Buzzer module}\quad S / VCC* / GND\\GPIO25 is configured OUTPUT, initially LOW.\\A fall-triggered backup alarm is not implemented.};
\draw[draw=warning,line width=1pt] (3.86,.8)--(4.1,.8)--(4.1,-1.55)--(8.1,-1.55);
\node[archlabel,text=warning] at (5.1,-1.3) {control S};
\draw[supply] (6,-1.9)--(8.1,-1.9);\fill[red!65!black] (6,-1.9) circle (.045);
\draw[ground] (6.65,-2.3)--(8.1,-2.3);\fill[ink] (6.65,-2.3) circle (.045);
\node[archnote,text width=14.6cm,anchor=north west] at (0,-3.12)
 {\textbf{Electrical boundary.} I\textsuperscript{2}C runs at 100 kHz in the source. Use ESP32-compatible 3.3 V logic and shared ground. VCC* means the supply required by the actual breakout: confirm its regulator, pull-ups and carrier rail before wiring. The bare MAX30102 needs separate core/LED supplies; do not infer its breakout pin rating from this functional drawing. Buzzer driver/current requirements depend on the module.\\[4pt]
  Microphone GPIO34 and light-sensor GPIO35 are retired. INT, AD0 and auxiliary IMU pins are not read by this source. Crossings without a dot are separate nets; dots identify joins.};
\end{tikzpicture}}
\caption{Functional connection map of the user-specified current KEYES/ESP32 assembly, derived from source GPIOs and I\textsuperscript{2}C addresses. It records the intended wiring rather than certifying every cable in the earlier photograph. The photograph predates removal of the microphone and light sensor. The retired input references were removed and the source compiled without flashing.}\label{fig:wearable-wiring}
\end{figure}}

\newcommand{\WearablePurposeFigure}{%
\begin{figure}[H]\centering
\resizebox{\linewidth}{!}{%
\begin{tikzpicture}[font=\sffamily\small]
\node[archbox,text width=4.25cm,font=\sffamily\scriptsize,minimum height=2.7cm] (ppg) at (2.35,6)
 {\textbf{1 / Optical contact}\quad MAX30102\\Red/IR signals support pulse and SpO\textsubscript{2} estimates. Contact and validity gates prevent missing signal becoming a zero reading. Bench estimates are not clinical validation.};
\node[archbox,text width=4.25cm,font=\sffamily\scriptsize,minimum height=2.7cm] (temp) at (12.65,6)
 {\textbf{2 / Thermal context}\quad GY906\\Object surface and local ambient are shown separately. Context motivates heat/cold research; surface temperature is not core body temperature.};
\node[archbox,text width=4.25cm,font=\sffamily\scriptsize,minimum height=2.7cm] (imu) at (2.35,2.1)
 {\textbf{3 / Movement input}\quad MPU6050\\Acceleration and angular velocity describe movement. Walking or falls require a separate time-window algorithm and validation; a one-second snapshot is insufficient.};
\node[archbox,text width=4.25cm,font=\sffamily\scriptsize,minimum height=2.7cm] (mobile) at (12.65,2.1)
 {\textbf{4 / Phone as gateway}\\Use its GPS, network, storage and notification speaker. Adult consent controls optional location/audio. Runtime depends on workload and batteries; it has not been measured.};
% Person, wrist module and phone are deliberately schematic.
\draw[ink,line width=1.4pt] (7.5,6.55) circle (.45);
\draw[ink,line width=2pt,line cap=round] (7.5,6.1)--(7.5,3.4)
 (7.5,5.5)--(6.2,4.35) (7.5,5.5)--(8.7,4.35)
 (7.5,3.4)--(6.55,1.55) (7.5,3.4)--(8.45,1.55);
\draw[accent,fill=accent!15,rounded corners=2pt,line width=1pt] (5.97,4.19) rectangle (6.5,4.58);
\draw[ink,fill=light,rounded corners=3pt] (8.52,3.65) rectangle (9.12,4.62);
\draw[ink!60] (8.7,3.8)--(8.93,3.8);
\node[archnote,align=center,text width=3.7cm] at (7.5,.72) {Older adult wearing the assembly;\\the illustration is a design scenario,\\not a participant photograph.};
\draw[archflow] (ppg.east)--(5.1,6)--(6.23,4.58);
\draw[archflow] (temp.west)--(9.9,6)--(6.5,4.5);
\draw[archflow] (imu.east)--(5.2,2.1)--(6.1,4.19);
\draw[archflow] (mobile.west)--(9.5,2.1)--(8.98,3.65);
\node[archfuture,text width=14.1cm,anchor=north] at (7.5,-.15)
 {\textbf{Future local fallback / GPIO25 buzzer.} An autonomous alarm during phone failure needs local detection logic and an independently powered wearable. The present source holds the output LOW. No current runtime or battery test establishes that fallback.};
\node[archnote,text width=14.5cm,anchor=north] at (7.5,-1.75)
 {\textbf{Interpretation.} No optical contact means no reliable optical measurement; it does not prove the wearable is off the body. Failed wiring, motion, poor placement or a disconnected sensor can also remove readings. These modules do not measure hydration: an agreed drinking reminder is a different feature.};
\end{tikzpicture}}
\caption{Why each component belongs in the design. Solid teal callouts describe available acquisition or phone responsibilities; the amber panel identifies the requested future backup behaviour. Thermal rules and machine-learning exercises remain educational and are not applied as diagnoses to physical wearable readings.}\label{fig:wearable-purpose}
\end{figure}}
